\documentclass[11pt,fleqn]{article} 
\usepackage[margin=0.8in, head=1in]{geometry} 
\usepackage{amsmath, amssymb, amsthm}
\usepackage{fancyhdr} 
\usepackage{palatino, url, multicol}
\usepackage{graphicx} 
\usepackage[all]{xy}
\usepackage{polynom} 
\usepackage{pdfsync}
\usepackage{enumerate}
\usepackage{framed}
\usepackage{setspace}
\usepackage{array,tikz}
\pagestyle{fancy} 
\lhead{Linear Algebra}
\chead{Midterm II Review}
%\rhead{\Large{Name: \underline{\hspace{2in}}}}

\newcommand{\vtr}[1]{\textbf{#1}}

\def\vectwo#1#2{\begin{bmatrix}#1\\#2\end{bmatrix}}
\def\vecthree#1#2#3{\begin{bmatrix}#1\\#2\\#3\end{bmatrix}}
\def\vecfour#1#2#3#4{\begin{bmatrix}#1\\#2\\#3\\#4\end{bmatrix}}
\def\bbm{\begin{bmatrix}}
\def\ebm{\end{bmatrix}}

% macros
\usepackage{amssymb}
\newcommand{\bA}{\mathbf{A}}
\newcommand{\bB}{\mathbf{B}}
\newcommand{\bE}{\mathbf{E}}
\newcommand{\bF}{\mathbf{F}}
\newcommand{\bQ}{\mathbf{Q}}
\newcommand{\bP}{\mathbf{P}}

\newcommand{\ba}{\mathbf{a}}
\newcommand{\bb}{\mathbf{b}}
\newcommand{\bc}{\mathbf{c}}
\newcommand{\br}{\mathbf{r}}
\newcommand{\bv}{\mathbf{v}}
\newcommand{\bw}{\mathbf{w}}
\newcommand{\bx}{\mathbf{x}}
\newcommand{\be}{\mathbf{e}}
\newcommand{\bp}{\mathbf{p}}
\newcommand{\bq}{\mathbf{q}}

\begin{document}
\renewcommand{\headrulewidth}{0pt}
\newcommand{\blank}[1]{\rule{#1}{0.75pt}}
\renewcommand{\d}{\displaystyle}

The exam will be Wednesday November 12, 9:15-10:15. No books, notes, calculator, electronics, or internet access.\\

It will cover Chapters 3 and 4 and Section 5.1 on Determinants.\\

\noindent Chapter 3: Vector Spaces and Subspaces\\
Section 3.1: Spaces of Vectors\\

terminology/definitions: \\
a \emph{vector space} and a \emph{subspace} of a vector space, \\
$\mathbb{R}^n$ as a vector space\\
the vector spaces $\mathbf{M, F, Z}$ (real 2 by 2 matrices, real functions $f(x)$ and the zero vector), \\
defining a subspace as the set of all linear combinations of a set of vectors\\
\emph{column space} of a matrix, $C(A)$\\
a subspace \emph{spanned} by a set of vectors, \\

skills/principles: \\
identify sets that are or are not subspaces of standard vector spaces, \\
the geometric properties of subspaces (when compared to sets that are not subspaces), \\
the equivalence of the statements $Ax=b$ is solvable and $b$ is in the column space of $A.$\\

\vskip 0.1in

\noindent Section 3.2: The Nullspace of $A$: Solving $Ax=0$ and $Rx=0$\\

terminology/definitions: \\
the \emph{null space of a matrix $A$}, $N(A)$\\
\emph{reduced row echelon form} of a matrix, typically $R$ (and this is different from $U$, notation for the upper triangular form)\\
\emph{pivot} and \emph{free} columns\\
\emph{special solution} associated with a free column\\
\emph{rank} of a matrix and what it indicates about a matrix, solutions to $Ax=b$, and free columns\\

skills/principles: \\
know how to find a reduced row echelon form of a matrix\\
know how to efficiently find the null space of a matrix\\
know that row operations do not change the null space\\

\noindent Section 3.3: The Complete Solutions to $Ax=b$ \\

terminology/definitions: \\
a \emph{complete} solution to $Ax=b$ and how to find it efficiently,\\
\emph{full column rank} or \emph{full row rank} and what these indicate about $N(A)$ or $C(A)$

skills/principles: \\
At the end of this section, we are supposed to understand in great detail what the reduced row echelon form of $A$ should indicate about the nature of solutions to $Ax=b,$ $C(A)$, and $N(A).$

\vskip 0.1 in

\noindent Section 3.4: Independence, Basis, and Dimension \\

terminology/definitions: \\
\emph{linearly independent} sets of vectors and \emph{linearly dependent} sets of vectors\\
\emph{a spanning set of vectors}\\
\emph{basis}\\
\emph{row space of matrix $A$}, $C(A^T)$\\
\emph{dimension} of a vector space\\

skills/principles: \\
how to show (rigorously) that at set of vectors is linearly independent or linearly dependent\\
how to show (rigorously) that a set of vectors spans a vector space (or subspace)\\
how to find bases for subspaces, like $C(A)$ or $C(A^T)$ or $N(A)$\\
how to determine the dimension of subspaces\\

\vskip 0.1 in

\noindent Section 3.5: Dimensions of the Four Subspaces\\

terminology/definitions: \\
the \emph{four subspaces} in words: \emph{column space, row space, null space} and \emph{left null space} and in symbols: $C(A), C(A^T), N(A), N(A^T)$\\


skills/principles: \\
know how to find each of the four subspaces, find bases for each of them, determine their dimension, and understand their geometry (what space to they live in? orthogonality?)\\
how elimination does (or does not) change the four subspaces\\
know the basic principles on page 184 of your text. (You don't need to be able to state the text's "Fundamental Theorem of Linear Algebra", but you should definitely know it!)\\

\vskip 0.1 in

\noindent Chapter 4: Orthogonality\\
\noindent Section 4.1: Orthogonality of the Four Subspaces \\

terminology/definitions: \\
\emph{orthogonal vectors} and the algebraic consequences\\
\emph{orthogonal vector spaces}, and examples\\
\emph{orthogonal complements} \\

skills/principles: \\
If one knows a vector space, $V$, has dimension $d$ and one has a set $S$ of $d$ vectors in $V$, one can know that $S$ is a basis of $V$ by checking \emph{one} of the two properties of a basis (ie check $S$ is linearly independent  OR check $S$ spans $V.$) One does not have to check both.


\vskip 0.1 in

\noindent Section 4.2: Projections \\


terminology/definitions: \\
the \emph{projection}, $\bp$, of a vector $\bb$ onto a vector $\ba$\\
the \emph{projection}, $\bp$, of a vector $\bb$ onto a subspace\\
the \emph{error} $\be=\bb-\bp$\\
$\mathbf{P}$, the projection matrix\\
$\widehat{x}$, the coefficient vector used to construct the projection $\bp$


skills/principles: \\
know how to compute the projection of one vector onto another and how to draw the projection, and the same for the error\\
know how to compute the projection of one vector onto a subspace and how to draw the projection, and the same for the error\\
know that the projection minimizes the error as measured by the sum of the squares of the differences in components\\
know how to compute $\mathbf{P}$, its properties, and how to use it\\
geometrically, we think of $\bp$ as the vector in $S$ closest to the vector $\bb$\\

\vskip 0.1 in

\noindent Section 4.3: Least Squares Approximations \\


terminology/definitions: \\
the \emph{normal equations}: $\bA^T\bA\bx=\bA^T\bb$, whose solution is $\widehat\bx$\\

skills/principles: \\
the projection $\bp$ (from 4.2) and the methods to find it can be reinterpreted to finding best solutions to over-determined systems of equations or to finding curves of best fit to a given data set.\\
when $\bA\bx=\bb$ has no solution, $\widehat\bx$, is the ``best" solution, or the ``least squares" solution because it minimizes the error $\lvert \bA\bx-\bb \rvert^2$\\
the least squares solutions can be used to find the line (or parabola or cubic...) of best fit to a set of points (data)\\

\vskip 0.1 in

\noindent Section 4.4: Orthonormal Bases and Gram-Schmidt\\


terminology/definitions: \\
\emph{orthogonal basis}\\
\emph{orthonormal basis}\\
\emph{Gram-Schmidt Process}

skills/principles: \\
how to find the projection of a vector $\bb$ onto a subspace $S$ given an orthonormal basis $\bq_1, \cdots \bq_n$\\
how to use the Gram-Schmidt Process to find an orthogonal or orthonormal basis from a given basis\\
If $\bQ$ has orthonormal columns, then know the many efficiencies that result. (ex: $\bQ^T\bQ=$?, finding $\widehat{\bx}$, finding $\bP$)

\vskip 0.1 in

\noindent Chapter 5 Section 1: The Properties of Determinants\\

terminology/definitions: \\
\emph{determinant} of a 2 by 2 matrix

skills/principles: \\
There are 10 principles to know about the determinant (listed in the text)\\
Know how to apply these to an $n$ by $n$ matri


\end{document}